% -----------------------------------------------------------------------------
% This file is automatically generated.
% Do not edit manually.
% Source: notebooks/support/regression-analysis/support.Rmd
% -----------------------------------------------------------------------------


\noindent In contrast, the stepwise-selection helper reports AIC values based on a
deviance-like approximation, that omits constant terms unrelated to model
comparison:
$\text{AIC}_{\text{step}} = n \log\left(\frac{\text{SSE}}{n}\right) + 2k$.
The difference between the two expressions is a constant shift of
$n \cdot (\log(2\pi) + 1) + 2$.
The term $n \cdot (\log(2\pi) + 1)$ reflects the part of the log-likelihood
that depends only on the sample size and the assumption of normally distributed
errors. The additional +2 arises because the full likelihood-based AIC
evaluates the likelihood based on a model with an estimated variance parameter
$\sigma^2$, introducing one more estimated parameter than is typically counted
in the model formula (e.g., intercept and slopes). This shifts the effective
count of parameters k by 1 and contributes an extra penalty of 2 in the AIC
expression.
\noindent Similarly, the BIC values are obtained as follows.

\par\addvspace{\topsep}
\begingroup
\fvset{listparameters={%
  \setlength{\topsep}{0pt}%
  \setlength{\partopsep}{0pt}%
  \setlength{\parsep}{0pt}%
  \setlength{\itemsep}{0pt}%
}}
\begin{Verbatim}[breaklines=true,formatcom=\color{ada_blue}]
print('BIC ussteamco_mod_0:', ussteamco_mod_0.bic)
print('BIC model_backward:', model_backward.bic)
print('BIC ussteamco_mod_1:', ussteamco_mod_1.bic)
print('BIC ussteamco_mod_2:', ussteamco_mod_2.bic)
\end{Verbatim}
\begin{Verbatim}[breaklines=true,formatcom=\color{ada_light_blue}]
[1] 1207.383

[1] 1156.611

[1] 1159.766

[1] 1157.468
\end{Verbatim}
\endgroup
\par\addvspace{\topsep}
\noindent The same consideration applies to the BIC. The BIC values of \ttblue{mod.0}, \ttblue{model\_backward}, \ttblue{mod.1}, and \ttblue{mod.2} may be compared directly because they are based on the same response data. The BIC of \ttblue{mod.3} should not be interpreted as evidence that the winsorized model is better or worse than the previous models, because it is based on a different dataset.
